\documentclass[12pt]{article}
\usepackage[utf8]{inputenc}
\usepackage[T1]{fontenc}
\usepackage[margin=1in]{geometry}
\usepackage{setspace}
\usepackage{xcolor}
\usepackage[colorlinks=true, urlcolor=blue!60!black, linkcolor=blue!60!black]{hyperref}
\setlength{\parindent}{2em}
\setlength{\parskip}{1.25em}
\renewcommand{\baselinestretch}{1.0}
\begin{document}
\begin{center}
{\LARGE \textit{The Slaughter-House Cases}, 83 U.S. 36 (1873)\par}
\vspace{0.8em}
{\large Case Brief --- Four-Page Edition\par}
\end{center}
\vspace{1em}

\section{Memory Jogger}

In its first interpretation of the Reconstruction Amendments, the Supreme Court upheld a 25-year Louisiana slaughterhouse monopoly and narrowly confined the Fourteenth Amendment's Privileges or Immunities Clause to rights owing their existence to the federal government rather than ordinary civil or economic rights of state citizenship (\href{https://www.courtlistener.com/opinion/88661/butchers-benevolent-assn-v-crescent-city-live-stock-landing/}{83 U.S. 36}).

\section{Detailed Case Facts}

In 1869, responding to slaughterhouse offal and livestock waste fouling the Mississippi River upstream of New Orleans's water supply, the Louisiana legislature barred landing, keeping, or slaughtering animals anywhere in the parishes of Orleans, Jefferson, and St.\ Bernard---1,154 square miles---except at designated facilities. It chartered the Crescent City Live-Stock Landing and Slaughter-House Company with a 25-year exclusive privilege, required it to build a ``grand slaughter-house,'' and ordered competitors closed, while obliging the company to let any butcher slaughter there at capped fees, subject to inspection.

Roughly 1,000 butchers had operated independent establishments. Represented by former Justice John A. Campbell, they challenged the monopoly as ``involuntary servitude'' under the Thirteenth Amendment and as abridging their privileges or immunities, depriving them of liberty and property without due process, and denying them equal protection under the Fourteenth.

\section{Procedural History}

Hundreds of suits were filed in Louisiana courts, including actions by the Butchers' Benevolent Association and an enforcement suit by the state Attorney General. The Louisiana Supreme Court upheld the statute. Three consolidated cases reached the U.S. Supreme Court on writs of error; after argument in 1872 and reargument before a full bench in February 1873, the Court affirmed on April 14, 1873.

\section{Judicial Votes}

Affirmed, \textbf{5--4}.
\begin{itemize}
\item \textbf{Majority:} Justice Miller, joined by Justices Clifford, Davis, Strong, and Hunt.
\item \textbf{Dissents:} Justice Field, joined by Chief Justice Chase and Justices Swayne and Bradley; Justice Bradley, separately; Justice Swayne, separately.
\end{itemize}

\section{Holding}

The Citizenship Clause creates two distinct citizenships, and the Privileges or Immunities Clause protects only privileges owing their existence to the federal government, its national character, Constitution, or laws, leaving fundamental civil and economic rights such as pursuing a trade to state police power. The Thirteenth Amendment bars personal servitude, not commercial monopolies, and the Due Process and Equal Protection Clauses do not invalidate an ordinary health-and-safety monopoly. Miller's list of national privileges and his doubt that the Equal Protection Clause reaches anything beyond discrimination against ``the negroes as a class'' are dicta.

\section{Analysis of Opinions}

\textbf{Majority (Miller).} The statute was a classic police-power regulation of an unwholesome trade, and it did not strip butchers of their trade because the company had to admit them all. The Amendments shared ``one pervading purpose'': the freedom and civil equality of the formerly enslaved. Because Section 1 distinguishes ``citizens of the United States'' from citizens ``of the State wherein they reside'' and protects only the former's privileges, the framers left state-citizenship rights---property and livelihood, as in \textit{Corfield v. Coryell}---to the states and Article IV. A broader reading would ``fetter and degrade the State governments'' and make the Court a ``perpetual censor upon all legislation of the States.''

\textbf{Field (dissenting).} Separating genuine sanitary provisions from the monopoly, Field called the health justification a ``shallow'' pretence. National citizenship is now primary, and the Clause protects \textit{Corfield}'s fundamental rights, above all ``the right to pursue a lawful employment in a lawful manner, without other restraint than such as equally affects all persons.'' Drawing on \textit{The Case of Monopolies} and the Statute of Monopolies, he argued the majority made the Clause ``a vain and idle enactment, which accomplished nothing.''

\textbf{Bradley (dissenting).} National citizenship carries both the rights of Englishmen, including choosing a calling, and the enumerated Bill of Rights guarantees. A monopoly excluding a large class from a lawful occupation violates due process (``Their right of choice is a portion of their liberty; their occupation is their property'') and equal protection; the federalism concern is an ``argument from inconvenience.''

\textbf{Swayne (dissenting).} The Amendments are a ``new Magna Charta'' standing ``at the opposite pole from the first eleven''; ``[l]abor is property,'' and ``the novelty was known and the measure deliberately adopted.''

\section{Examples: Future Applications}

\begin{enumerate}
\item \textbf{Same-Side (Upheld): Waste-Hauling Franchise.} A state limits waste hauling to one regulated contractor serving all at capped rates. A displaced hauler's Privileges or Immunities claim fails: pursuing a local trade is a state-citizenship right.
\item \textbf{Same-Side (Upheld): State Bar Admission.} As \textit{Bradwell v. Illinois} held the next day, practicing law in state courts is not a privilege of national citizenship.
\item \textbf{Opposite-Side (Invalid): Travel Penalties.} A head tax on travel to Washington, D.C., or penalties for moving to another state abridge privileges Miller himself recognized.
\item \textbf{Opposite-Side (Invalid): Federal Waterways.} Barring federally licensed vessels from navigable waters interferes with rights that ``owe their existence to the Federal government.''
\item \textbf{Fence-Sitter: Incorporation.} A challenge to an arms or excessive-fines law under the Privileges or Immunities Clause. \textit{McDonald} retained due process incorporation, but Miller's mention of assembly and habeas corpus fuels calls by Justices Thomas and Gorsuch to revisit the case.
\end{enumerate}

\section{Critique}

\textbf{Originalist and textualist views.} Originalists widely see the case as a grave textual error that erased the Clause. Randy Barnett and Evan Bernick argue it protected both enumerated and unenumerated rights, including an honest livelihood, as Senator Jacob Howard explained; Kurt Lash reads it as limited to enumerated rights but still rejects Miller. Justice Thomas urges overruling the case.

\textbf{Progressive and doctrinal views.} Many condemn the case as a retreat from Reconstruction that forced incorporation through due process, yet defend Miller's result upholding a biracial legislature's health law and note that Field's and Bradley's dissents seeded \textit{Lochner}. Bryan Wildenthal and Kevin Newsom argue \textit{Cruikshank} bears more blame.

\section{Key Quotations}

\begin{quote}
``We repeat, then, in the light of this recapitulation of events, almost too recent to be called history, but which are familiar to us all; and on the most casual examination of the language of these amendments, no one can fail to be impressed with the one pervading purpose found in them all, lying at the foundation of each, and without which none of them would have been even suggested; we mean the freedom of the slave race, the security and firm establishment of that freedom, and the protection of the newly-made freeman and citizen from the oppressions of those who had formerly exercised unlimited dominion over him.''
\end{quote}
\noindent --- Justice Miller, for the Court, 83 U.S. at 71.

\begin{quote}
``And still further, such a construction followed by the reversal of the judgments of the Supreme Court of Louisiana in these cases, would constitute this court a perpetual censor upon all legislation of the States, on the civil rights of their own citizens, with authority to nullify such as it did not approve as consistent with those rights, as they existed at the time of the adoption of this amendment.''
\end{quote}
\noindent --- Justice Miller, for the Court, 83 U.S. at 78.

\begin{quote}
``If this inhibition has no reference to privileges and immunities of this character, but only refers, as held by the majority of the court in their opinion, to such privileges and immunities as were before its adoption specially designated in the Constitution or necessarily implied as belonging to citizens of the United States, it was a vain and idle enactment, which accomplished nothing, and most unnecessarily excited Congress and the people on its passage.''
\end{quote}
\noindent --- Justice Field, dissenting, 83 U.S. at 96.

\end{document}
